\section{Four complete exceptional small-field core exclusions}
\label{sec:exceptional-small-cores}

\paragraph{Status.}
New complete root derivation, read in full by the author and an
independent reviewer with both mathematical audits passing. Full GAP
orthogonality and spectrum checks, a separate root exact reconstruction,
and an independent polynomial-Fourier reconstruction all passed.
The original character tables are cited inputs with their distinct
provenances; they have not been independently constructed in this run.

\begin{lemma}[D-EXS-01: complete core and outer spectra]
\label{D-EXS-01}
Let $W$ be a nontrivial ordinary irreducible complex representation of a
quasisimple group with one of the simple quotients in the first column
below. Put $a=\dim W$ and
$\nu(T)=a-\max_\lambda\dim\ker(T-\lambda I)$. Every nonscalar core
operator satisfies the indicated bound:
\[
\begin{array}{c|c|c|c}
\text{simple quotient}&\text{universal cover used}&a\text{ at least}
                                      &\nu(\rho(x))/a\text{ at least}\\\hline
F_4(2)&2.F_4(2)&52&4/13\\
F_4(3)&F_4(3)&6643&32/73\\
E_6(2)&E_6(2)&2482&512/1241\\
{}^2E_6(2)&(2^2\times3).{}^2E_6(2)&1938&8/19
\end{array}                                                    \tag{EXS.1}
\]
For the first row, every nonscalar finite-order operator normalizing
the represented image has the same bound $\nu(T)\ge4a/13$.
No faithfulness on the entire center is assumed.
\end{lemma}

\begin{proof}
The universal-cover classification is used in precisely the ordinary
complex sense. The exceptional multiplier cases here are $F_4(2)$ and
${}^2E_6(2)$; the complete list is
\cite[Chapter~5, Table~1, p.~35]{BI:HHM}. The exact CTblLib metadata
\texttt{ExtensionInfoCharacterTable} for the four simple quotients,
in the order of (EXS.1), are
\[
 [\texttt{2},\texttt{2}],\quad
 [\texttt{empty},\texttt{empty}],\quad
 [\texttt{empty},\texttt{2}],\quad
 [\texttt{(2\^{}2x3)},\texttt{3.2}].                         \tag{EXS.2}
\]
The first entry specifies the multiplier and the second the outer
automorphism group. In particular the last center is noncyclic.
We retain \emph{every} nonlinear row of its universal table, so all
central quotients and all possible central-character kernels are included.
We do not replace this table by a presumed faithful representation
of the full noncyclic center.

The exact source identifiers and provenance are in
\cite{HD:ExceptionalSmallTables}. The $F_4(2)$ cover and its outer
extension are ATLAS inputs; the retained extension uses the corrected
square map. The $F_4(3)$ table was computed by Eamonn O'Brien with
Magma in May/June~2024 and added to CTblLib~1.3.10. The $E_6(2)$
table is Fischer's table with Unger's 2016 corrections to order-91
power maps and irrationalities. The universal ${}^2E_6(2)$ table
records its construction by \texttt{PossibleCharacterTablesOfTypeV4G}.
These separate data origins are not all relabelled as printed ATLAS
tables. The actual export uses CTblLib~1.3.11.

The complete mathematical input is
\texttt{research/exceptional\_small\_inputs.json}, with SHA-256
\texttt{625618617ad0b6268b8898cc39e7e04afb6ae55f8639686347a0b9cf8d244a21}.
It retains every class size, element order, ordinary character value,
prime power map, indicator, source string, and the core-to-outer fusion.
For an element $x$ of order $d$, every eigenvalue multiplicity is
computed by the exact identity
\[
 m_j=\frac1d\sum_{u=0}^{d-1}\chi(x^u)\exp(-2\pi iju/d),
                       \qquad 0\le j<d.                        \tag{EXS.3}
\]
The reconstruction uses rational field traces and Ramanujan sums,
with exact integer arithmetic, and verifies integrality,
nonnegativity and the sum $\sum m_j=a$. It also checks the degree
sum, all power-map orders, and the normal fusion's identity kernel,
complete class unions and compatibility with prime power maps.
An entry is skipped only when $|\chi(x)|^2=a^2$, which for a
finite-order operator is exactly the scalar condition.

For the table \texttt{2.F4(2).2}, a row is selected exactly when
its restriction to the core has norm one and matches an actual
nontrivial irreducible core row. The complete counts and minima are
\[
\begin{array}{c|r|r|c}
\text{table}&\text{selected rows}&\text{nonscalar spectra}&\min\nu/a\\\hline
\texttt{2.F4(2)}&169&28392&4/13\\
\texttt{2.F4(2).2}&86&12814&4/13\\
\texttt{F4(3)}&272&73984&32/73\\
\texttt{E6(2)}&179&32041&512/1241\\
\texttt{(2\^{}2x3).2E6(2)}&933&860226&8/19
\end{array}                                                    \tag{EXS.4}
\]
Thus there are $1639$ selected rows and $1007457$ spectra.
The script \texttt{check\_exceptional\_small.py} retains the minimum
and a witnessing class for every row in
\texttt{exceptional\_small\_checked.json}. These are complete exact
finite inequalities, not numerical sampling or a calculation of only
unipotent characters. The arithmetic helper is the previously reviewed
Ramanujan implementation, reused here separately from GAP's
\texttt{EigenvaluesChar}; it is not advertised as a new third algorithm.

For completeness the outer-table assertion covers actual normalizers,
not only the chosen abstract extension. The universal cover of $F_4(2)$
has center of order two and outer group of order two. An automorphism
of a represented central quotient lifts uniquely to the universal
perfect cover. If it induces the identity on the simple quotient,
its difference from the identity is a homomorphism from a perfect
group to its center and is trivial. Inner automorphisms are implemented
by core operators up to scalar, by Schur's lemma.

For an outer automorphism, adjust by an inner operator to the
automorphism induced by a fixed lift $t$ in \texttt{2.F4(2).2}.
A compatible intertwiner $A$ satisfies $A^2=c\rho(t^2)$ for a
nonzero scalar $c$. Rescale $A$ to remove $c$. The defining relation
then gives a genuine representation of this extension with irreducible
core restriction, exactly among the selected rows. The original
normalizer is a scalar multiple of an element of that representation.
This does not assume a split extension and also handles a different
isoclinic scalar square. Since scalar multiplication preserves $\nu$,
(EXS.4) proves the asserted full-normalizer bound.
\end{proof}

\begin{lemma}[D-EXS-02: uniform support for all full normalizers]
\label{D-EXS-02}
For every representation in D-EXS-01, every nonscalar finite-order
operator normalizing the represented core has
\[
                              \nu(T)\ge a/5.                   \tag{EXS.5}
\]
This includes all compatible outer automorphisms, without assuming
selfduality or a cyclic outer group in the last row.
\end{lemma}

\begin{proof}
The first row already has the stronger bound $4a/13$.
In each of the other three rows, (EXS.1) gives core support
at least $2a/5$: indeed $5\cdot32\ge2\cdot73$,
$5\cdot512\ge2\cdot1241$, and $5\cdot8\ge2\cdot19$.

Choose a represented core element $h$ with $[T,h]\ne I$;
otherwise irreducibility and Schur's lemma would make $T$ scalar.
The commutator belongs to the represented core. If it is nonscalar,
its moving rank is at least its projective support, hence at least
$2a/5$. If it is a nonidentity scalar, its moving rank is $a$.
Normalize $T$ by a scalar so that a largest eigenspace has eigenvalue
one. This does not change its commutator, and
\[
 \operatorname{rank}([T,h]-I)
   =\operatorname{rank}(Th-hT)
   \le2\operatorname{rank}(T-I)=2\nu(T).                       \tag{EXS.6}
\]
The equality follows by multiplying by invertible matrices, and the
inequality by writing $Th-hT=(T-I)h-h(T-I)$. Thus (EXS.5) follows.
No description of the outer automorphism action on the noncyclic
center of the last universal cover is needed for this commutator
argument; every compatible normalizer is handled directly.
\end{proof}

\begin{theorem}[D-EXS-03: no mixed products of these exceptional cores]
\label{D-EXS-03}
In the actual sixfold curve-cover source problem, an inducing block
cannot have period rank nine if its irreducibly acting normal
quasisimple central-product core has all simple quotients among
\[
                  F_4(2),\quad F_4(3),\quad E_6(2),
                             \quad {}^2E_6(2).                \tag{EXS.7}
\]
Mixed types, arbitrary central quotients, compatible outer and scalar
operators, and permutations of factors are permitted. No transitivity,
selfduality, character-field equality or polarization determinant
hypothesis is imposed beyond the stated actual source setting.
\end{theorem}

\begin{proof}
Lift the core to the product of its universal covers and write
$V=\bigotimes_i W_i$, with $e=\prod_i a_i$. Every $a_i\ge52$.
For a factor-preserving nonscalar normalizer, at least one factor
operator is nonscalar. Normalized eigenvalue multiplicities convolve
under tensor product; their maximum is no greater than the maximum
for that factor. D-EXS-02 therefore gives support at least $e/5$.

For a nontrivial factor permutation, the complete weighted basis-word
argument of D-SPOR-03 gives the stronger $e/4$. Explicitly, on a
cycle of length $l\ge2$ with factor dimension $a\ge2$, diagonalize
the product of the intertwiners around the cycle. In the resulting
weighted word basis, at most $a$ of the $a^l$ words are fixed;
every other orbit has length at least two and contributes at most
one dimension to any chosen eigenspace. Thus its largest eigenspace
has dimension at most $(a^l+a)/2\le3a^l/4$. Tensoring the other
cycles preserves this upper fraction. This argument does not require
the intertwiners to be pure coordinate permutations. Finally a
nonidentity scalar inertia moves all $e$ dimensions.

The actual intermediate-cover Euler identity from D-BI-01 is
$6=2e(h-1)+D$, with $D$ the total inertia moving rank. Since
$e\ge52$, genus $h\ge2$ is impossible. At genus one, $D=6$
allows at most six genuine branches and actual period rank at most
six. Rank nine must consequently have $h=0$, at least twelve
genuine branch points, and $D=2e+6$. Applying (EXS.5) at each
branch gives
\[
                 2e+6\ge12e/5,\qquad e\le15,                 \tag{EXS.8}
\]
contrary to $e\ge52$. This proves the theorem for every allowed
mixture, without a reality or multiplier argument.
\end{proof}

\paragraph{Geometry, framework, and precise boundary.}
This is a restriction on actual finite covers and their weight-one
multiplicity variations. Its spectrum and tensor identities lie in
the representation ring; the scheme-spreading, SGA finite-cover,
lambda-ring, K-theory and etale-weight placements have the exact
locators in the framework section. It strictly narrows the supplier
gate for condition A, cote16 scan050. It produces no new effective
cycle or detector constant, and changes neither condition I at
cote16 scan049 and cote44 Theorem4 nor the Cayley--Hamilton projector
at cote16 scans048--049. No motivic $L$-value estimate or general
negative Hodge-conjecture conclusion is inferred. Larger fields and
other exceptional families are outside (EXS.7).
